\section{Gaussian Sums}\label{sec:gaussian-sums}
Having studied additive and multiplicative characters of $\bbF_q$ separately, we now turn to hybrid sums that interleave both types. The first and most fundamental such object is the \emph{Gaussian sum}.\index{Gaussian sum}\index{Gauss, Carl Friedrich} A Gaussian sum is formed by taking the product of a multiplicative character and an additive character and summing over all nonzero elements of $\bbF_q$. The interplay between the two kinds of characters gives remarkable algebraic and analytic properties, and they will serve as the key tool throughout this chapter. We define them as follows.
\begin{Definition}{Gaussian Sum}{def:gaussian-sum}
  \index{Gaussian sum}%
  Let $\psi $ be a multiplicative character and $\chi$ be an additive character of $\bbF_q$. Then the \emph{Gaussian Sum} associated to $\psi$ and $\chi$ is defined as $$G(\psi,\chi)=\sum_{\alpha\in\bbF_q^*}\psi(\alpha)\cdot \chi(\alpha)$$
\end{Definition}

The absolute value of the Gaussian sum is at most $q-1$ by the triangle inequality\index{triangle inequality bound}. However, we will see below that when both $\psi$ and $\chi$ are non-trivial, $|G(\psi,\chi)|=\sqrt{q}$\index{Gaussian sum!absolute value formula} --- far smaller than this naive bound.

\begin{Theorem}{Evaluation of Gaussian Sums}{thm:gaussian-sum}
  Let $\psi$ be a multiplicative character and $\chi$ be an additive character of $\bbF_q$. Then \begin{equation}
    G(\psi,\chi)=\begin{cases}
      q-1 & \text{if $\psi=\psi_0$, $\chi=\chi_0$}     \\
      -1  & \text{if $\psi=\psi_0$, $\chi\neq \chi_0$} \\
      0   & \text{if $\psi\neq \psi_0$, $\chi=\chi_0$}
    \end{cases}
  \end{equation}
  If $\psi$ is non-trivial multiplicative character and $\chi$ is non-trivial additive character then $|G(\psi,\chi)|=\sqrt{q}$.
\end{Theorem}
\begin{proof}
  If $\psi=\psi_0$ and $\chi=\chi_0$ then for all $\alpha\in\bbF_q^*$ we have $\psi(\alpha)=1$ and $\chi(\alpha)=1$. Hence $G(\psi_0,\chi_0)=q-1$.

  Now suppose $\psi=\psi_0$ and $\chi$ is non-trivial additive character. Then by \thmref[(i)]{thm:add-char-props} $$G(\psi_0,\chi)=\sum_{\alpha\in\bbF_q^*}\chi(\alpha)=\sum_{\alpha\in\bbF_q}\chi(\alpha)-\chi(0)=-1$$

  Let $\psi$ be a non-trivial multiplicative character and $\chi=\chi_0$. Then for all $\alpha\in\bbF_q^*$ we have $\chi(\alpha)=1$. Hence by \thmref[(i)]{thm:mult-char-props}$$G(\psi,\chi_0)=\sum_{\alpha\in\bbF_q^*}\psi(\alpha)=0$$

  So now assume both $\psi$ and $\chi$ are non-trivial multiplicative and additive characters of $\bbF_q$ respectively. Then we have \begin{align*}
    |G(\psi,\chi)|^2 & = \ov{G(\psi,\chi)}\cdot G(\psi,\chi)                                                                                                                                                    \\
                     & = \lt(\sum_{\alpha\in\bbF_q^*}\ov{\psi(\alpha)}\cdot \ov{\chi(\alpha)}\rt) \cdot \lt(\sum_{\beta\in\bbF_q^*}\psi(\beta)\cdot \chi(\beta) \rt)                                            \\
                     & = \sum_{\alpha\in\bbF_q^*}\sum_{\beta\in\bbF_q^*}\ov{\psi(\alpha)}\cdot \ov{\chi(\alpha)}\cdot \psi(\beta)\cdot \chi(\beta)                                                              \\
                     & = \sum_{\alpha\in\bbF_q^*}\sum_{\beta\in\bbF_q^*}\psi(\beta\cdot \alpha^{-1})\cdot \chi(\beta-\alpha)                                                                                    \\
                     & = \sum_{\alpha\in\bbF_q^*}\sum_{\gm\in\bbF_q^*}\psi(\gm)\cdot \chi(\alpha(\gm-1))                                                             & [\text{Take }\gm=\beta\cdot \alpha^{-1}] \\
                     & = \sum_{\gm\in\bbF_q^*}\psi(\gm)\cdot \lt(\sum_{\alpha\in\bbF_q^*}\chi(\alpha(\gm-1))\rt)                                                                                                \\
                     & = \sum_{\gm\in\bbF_q^*}\psi(\gm)\cdot \lt(\sum_{\alpha\in\bbF_q}\chi(\alpha(\gm-1))\rt) -\chi(0)\sum_{\gm\in\bbF_q^*}\psi(\gm)                                                           \\
                     & = \sum_{\gm\in\bbF_q^*}\psi(\gm) \cdot \lt(\sum_{\alpha\in\bbF_q}\chi(\alpha(\gm-1))\rt)                                                      \by{\thmref[(i)]{thm:mult-char-props}}
  \end{align*}
  Hence by \thmref[(i)]{thm:add-char-props} $$\sum_{\alpha\in\bbF_q}\chi(\alpha(\gm-1))=\begin{cases}
      0 & \text{if $\gm\neq 1$} \\
      q & \text{if $\gm=1$}
    \end{cases}$$
  Therefore $|G(\psi,\chi)|^2=\psi(1)q=q$ and so we have the required result.
\end{proof}
Now under various transformations of the additive or multiplicative characters Gaussian sums gives some useful identities. We will see some of these below.

\begin{lemma}{Gaussian Sum Identities}{lem:gaussian-sum-identity}
  Gaussian sums for the finite field $\bbF_q$ satisfy the following identities. For any $a\in \bbF_q$ let $\chi_a$ be the additive character defined as in \thmref{thm:add-char-define}. Let $\psi\in\Mq$ be any multiplicative character and $\chi\in\Xq$ be any additive character of $\bbF_q$.  Then we have the following identities\begin{enumerate}[label=(\roman*)]
    \item $G(\psi,\chi_{ab})=\ov{\psi(a)}\cdot G(\psi,\chi_b)$ for any $a\in \bbF_q^*$ and $b\in \bbF_q$.
    \item $G(\psi,\ov{\chi})=\psi(-1)\cdot G(\psi,\chi)$.
    \item $G(\ov{\psi},\chi)=\psi(-1)\ov{G(\psi,\chi)}$.
    \item $G(\psi,\chi)\cdot G(\ov{\psi},\chi)=\psi(-1)q$ if $\psi$ and $\chi$ are non-trivial multiplicative and additive characters of $\bbF_q$ respectively.
    \item $G(\psi^p,\chi_b)=G(\psi,\chi_{\sg(b)} )$ for any $b\in \bbF_q$ where $p=\Char(\bbF_q)$ and $\sg(X)=X^p$ be the Frobenius map of $\bbF_q$.
  \end{enumerate}
\end{lemma}
\begin{proof}
  \begin{enumerate}[label=(\roman*)]
    \item For any $\gm\in\bbF_q$ we have $\chi_{ab}(\gm)=\chi_1(ab\cdot \gm)=\chi_b(a\cdot \gm)$. Hence $$G(\psi,\chi_{ab})=\sum_{\gm\in\bbF_q^*}\psi(\gm)\cdot \chi_{ab}(\gm)=\sum_{\gm\in\bbF_q^*}\psi(\gm)\cdot \chi_b(a\cdot \gm)$$Now take $d=a\cdot \gm$. Then $\gm=a^{-1}\cdot d$. Hence $$G(\psi,\chi_{ab})=\sum_{d\in\bbF_q^*}\psi(a^{-1}\cdot d)\cdot \chi_b(d)=\psi(a^{-1})\sum_{d\in\bbF_q^*}\psi(d)\cdot \chi_b(d)=\ov{\psi(a)}\cdot G(\psi,\chi_b)$$Hence we have the first identity.
    \item Now we have $\chi=\chi_b$ for some suitable $b\in \bbF_q$. Therefore $\ov{\chi}=\chi_{(-1)\cdot b}$. So by the first identity we have $G(\psi,\ov{\chi})=\ov{\psi(-1)}\cdot G(\psi,\chi_b)=\psi(-1)\cdot G(\psi,\chi)$. The last equality holds because $\psi(-1)=\pm 1$.
    \item Now again we have $\ov{\chi}=\chi_b$ for some suitable $b\in \bbF_q$. Hence $\chi=\chi_{(-1)\cdot b}$. Again by using the first identity we have $$G(\ov{\psi},\chi)=G\lt(\ov{\psi},\chi_{(-1)\cdot b}\rt)=\ov{\psi(-1)}\cdot G(\ov{\psi},\chi_b)=\psi(-1)\cdot G(\ov{\psi},\ov{\chi})$$Now from the definition of Gaussian sums we have $G(\ov{\psi},\ov{\chi})=\ov{G(\psi,\chi)}$. Hence we have the required identity.
    \item Using part (ii) and (iii) we have $$G(\psi,\chi)\cdot G(\ov{\psi},\chi)=G(\psi,\chi)\cdot \psi(-1)\ov{G(\psi,\chi)}=\psi(-1)\cdot |G(\psi,\chi)|^2$$Since $\psi$ and $\chi$ are non-trivial multiplicative and additive characters of $\bbF_q$ respectively by \thmref{thm:gaussian-sum} we have $|G(\psi,\chi)|^2=q$. Hence we have the required identity.
    \item If $\tr$ is the absolute trace function from $\bbF_q$ to $\bbF_p$ then for any $\alpha\in\bbF_q$ we have $\tr(\alpha)=\tr(\alpha^p)$. Hence $\chi_1(\alpha)=\chi_1(\alpha^p)$. Hence for any $\alpha\in \bbF_q$, we get $\chi_b(\alpha)=\chi_1(b\cdot \alpha)=\chi_1(b^p\cdot \alpha^p)=\chi_{\sg(b)}(\alpha^p)$. Therefore $$G(\psi^p,\chi_b)=\sum_{\alpha\in\bbF_q^*}\psi^p(\alpha)\cdot \chi_b(\alpha)=\sum_{\alpha\in\bbF_q^*}\psi(\alpha^p)\chi_{\sg(b)}(\alpha^p)=G(\psi,\chi_{\sg(b)})$$Therefore  we have the required result.
  \end{enumerate}
\end{proof}

Now we can use Gaussian Sums in a variety of context. For any multiplicative character $\psi$ we can express the value of $\psi$ at any nonzero $\alpha\in\bbF_q$ using Gaussian sums and additive characters. And similarly for any additive character $\chi$ of $\bbF_q$ we can express the value of $\chi$ at any nonzero $\alpha\in\bbF_q$ using Gaussian sums and multiplicative characters.

Now by \thmref[(iii)]{thm:add-char-props} for any $\alpha\in\bbF_q^*$ we have \begin{align*}
  \psi(\alpha) & =\frac1q\sum_{\beta\in \bbF_q^*}\psi(\beta)\sum_{c\in\bbF_q}\chi_{c}(\alpha)\ov{\chi_{c}(\beta)}  \\
               & =\frac1q\sum_{c\in\bbF_q}\chi_c(\alpha)\sum_{\beta\in\bbF_q^*}\psi(\beta)\cdot \ov{\chi}_c(\beta) \\
               & =\frac1q\sum_{c\in\bbF_q}\chi_c(\alpha)\cdot G(\psi,\ov{\chi}_c)                                  \\
               & =\frac1q\sum_{\chi\in\Xq}G(\psi,\ov{\chi})\chi(\alpha)
\end{align*}This may be thought of as the \emph{Fourier expansion}\index{Fourier expansion!of multiplicative characters}\index{Fourier expansion!multiplicative} of $\psi$ in terms of the additive characters of $\bbF_q$, with Gaussian sums appearing as Fourier coefficients. 

We can do similar thing for multiplicative characters of $\bbF_q$. By using \thmref[(iii)]{thm:mult-char-props} we get $$\chi(\alpha)=\frac1{q-1}\sum_{\psi\in\Mq}G(\ov{\psi},\chi)\chi(\alpha)$$ This can be interpreted as the Fourier expansion of the multiplicative characters of $\bbF_q$ with the Gaussian sums as Fourier coefficients. Hence we have the observations:
\begin{observation}\label{obs:fourier-coeff-gaussian-sum}
  For any multiplicative character $\psi$ and additive character $\chi$ of $\bbF_q$ and any  $\alpha\in\bbF_q^*$ we have \begin{enumerate}[label=(\roman*)]
    \item $\psi(\alpha)=\frac1q\sum\limits_{\chi\in \Xq}G(\psi,\ov{\chi})\chi(\alpha)$ 
    \item $\chi(\alpha)=\frac1{q-1}\sum\limits_{\psi\in\Mq}G(\ov{\psi},\chi)\psi(\alpha)$
  \end{enumerate}
\end{observation}


We now a another special formula for Gaussian sums for $2$-degree extension of finite fields and it  has a lot of application. But this has a restriction on the underlying field. 
\begin{Theorem}{Stickelberger's Theorem}{thm:stickelberger}%
  \index{Stickelberger's theorem}
  Let $q$ be a prime power and $\psi\in\Mq$ be a non-trivial multiplicative character of $\bbF_{q^2}$ of order $m$ such that $m\mid q-1$. Let $\chi_1$ be the canonical additive character of $\bbF_{q^2}$. Then $$G(\psi,\chi_1)=\begin{cases}
    q & \text{if $m$ is odd or $\frac{q+1}m$ is even}\\
    -q & \text{if $m$ is even and $\frac{q+1}m$ is odd}
  \end{cases}$$
\end{Theorem}
\begin{proof}
  For brevity of notations we write $E=\bbF_{q^2}$. Let $\gm$ be a primitive element of $E$. So set $g=\gm^{q+1}$. Therefore $g^{q-1}=1$ and so we get $g$ as primitive element of $\bbF_q$. Therefore for all $\alpha\in E^*$ there exists unique $j\in\{0,1,\dots,q-1$ and $k\in\{0,\dots, q+1\}$ such that $\alpha=g^j\cdot \gm^k$. Hence $\psi(g)=\psi^{q+1}(\gm)=1$ as $\ord(\psi)=m\mid q+1$. Therefore we have 
  \begin{align*}
    G(\psi,\chi_1) & = \sum_{j=0}^{q-2}\sum_{k=0}^{q}\psi(g^j\cdot \gm^k)\cdot \chi_1(g^j\cdot \gm^k)\\ 
    & = \sum_{k=0}^q\psi^k(\gm)\sum_{j=0}^{q-2}\chi_1(g^j\cdot \chi^k)\\ 
    & = \sum_{k=0}^q\psi^k(\gm)\sum_{\alpha\in\bbF_q^*}\chi_1(\alpha\cdot \chi^k)
  \end{align*} 
  Let $\tau_1$ is the canonical additive character of $\bbF_q$ then for any $\alpha\in\bbF_q^*$ and $k\in\{0,1,\dots,q\}$ we get $$\chi_1(\alpha\cdot \gm^k)=\tau_1\lt(\tr_{E/\bbF_q}(\alpha\cdot \gm^k)\rt)=\tau_1\lt(\alpha\cdot\tr_{E/\bbF_q}( \gm^k)\rt)$$Hence we get $$\sum_{\alpha\in\bbF_q^*}\chi_1(\alpha\cdot \chi^k)=\sum_{\alpha\in\bbF_q^*}\tau_1\lt(\alpha\cdot\tr_{E/\bbF_q}( \gm^k)\rt)=\begin{cases}
    -1 & \text{if }\tr_{E/\bbF_q}( \gm^k)\neq 0\\ 
    q-1 & \text{if }\tr_{E/\bbF_q}( \gm^k)=0
  \end{cases}$$Now since $[E:\bbF_q]=2$ we have $\tr_{E/\bbF_q}( \gm^k)=\gm^k+\gm^{k\cdot q}$. Therefore $$\tr_{E/\bbF_q}( \gm^k)=0\iff\gm^k=-\gm^{k\cdot q}\iff \gm^{k(q-1)}=-1$$So if $q$ is odd then $\gm^{k(q-1)}=-1\iff k=\frac{q+1}2$ and therefore we have $$\sum_{\alpha\in\bbF_q^*}\chi_1(\alpha\cdot \chi^k)=\begin{cases}
    -1 & \text{if }k\neq \frac{q+1}2\\
    q-1 & \text{if }k=\frac{q+1}2  
  \end{cases}$$And by using \thmref{thm:add-char-props} $$G(\psi,\chi_1) = q\cdot\psi^{\nicefrac{q+1}{2}}(\gm) - \sum_{k=0}^{q}\psi^k(\gm)=q\cdot\psi^{\nicefrac{(q+1)}{2}}$$
  Recall that $\psi$ has order $m$ and $m\mid q+1$. Therefore $\psi^{\nicefrac{(q+1)}{2}}(\gm)=\pm1$. Now if $\frac{q+1}{m}$ is even then we can rewrite $$\psi^{\nicefrac{(q+1)}{2}}(\gm)=\lt(\psi^m(\gm)\rt)^{\nicefrac{(q+1)}{2m}}=1$$ So suppose $\frac{q+1}{m}$ is odd, say $\frac{q+1}{m}=2\ell+1$ for some $\ell\geq 0$. Since $q$ is odd, $q+1$ is even, and as $(q+1)/m$ is odd the factor of $2$ in $q+1$ must lie entirely in $m$; in particular $m$ is even. Writing $\frac{q+1}{2}=m\ell+\frac{m}{2}$, we get $$\psi^{\nicefrac{(q+1)}{2}}(\gm)=\lt(\psi^m(\gm)\rt)^{\ell}\cdot \psi^{\nicefrac{m}{2}}(\gm)=\psi^{\nicefrac{m}{2}}(\gm).$$ Now $\psi^{\nicefrac{m}{2}}(\gm)$ is a square root of unity since order of $\psi=m$ and $\gm$ is primitive element of $E$ we have $\psi^{\nicefrac{(q+1)}{2}}(\gm)=-1$. Therefore $G(\psi,\chi_1)=(-1)^{\nicefrac{(q+1)}{m}}\cdot q$. So we have the theorem. 
\end{proof}

We can give a relation between Gaussian sums a multiplicative and additive character and Gaussian sum of lifted characters. To build that relation we need the theory of $L$-functions in \autoref{sec:l-functions}. Later we proved the \DHthm which gives this relation. 

Before going into this in the next section we will define another hybrid sum called \emph{Jacobi Sums}. Like Gaussian sum this also has many important applications. In the later sections we will use both Gaussian sums and Jacobi Sums to prove many interesting results